\documentclass[12pt,journal,compsoc]{IEEEtran}

\usepackage[english,portuguese]{babel}
\usepackage[utf8]{inputenc}
\usepackage[T1]{fontenc}
\usepackage[cmex10]{amsmath}
\usepackage{amssymb}
\usepackage{graphicx}
\usepackage{tikz}
\usetikzlibrary{arrows,positioning}
\usepackage{float}
\usepackage{xcolor}
\usepackage{listings}
\usepackage{hyperref}
\usepackage{fancyhdr}
\pagestyle{fancyplain}
\setlength{\headheight}{15pt}
\fancyhf{}
\lhead{1ELE910 -- Microcontroladores II}
\rhead{Roteiro: áudio no STM32F4-Discovery}
\cfoot{\thepage}
\hypersetup{colorlinks=true,linkcolor=blue,citecolor=blue,urlcolor=blue}
\lstset{basicstyle=\ttfamily\footnotesize,breaklines=true,frame=single,showstringspaces=false}
\newcommand{\funcao}[1]{\path{#1}}

\begin{document}
\title{Reprodução de Áudio no STM32F4-Discovery\\Roteiro Técnico Experimental}
\author{Projeto de Microcontroladores II}
\markboth{1ELE910 -- Microcontroladores II}{Reprodução de áudio no STM32F4-Discovery}
\IEEEtitleabstractindextext{%
\begin{abstract}
Este roteiro orienta a reprodução de áudio na STM32F4-Discovery, primeiro pela Flash interna e depois por um pendrive USB. A atividade parte do exemplo Audio playback and recording da STMicroelectronics, incluído no firmware STM32F4-Discovery V1.1.0. Siga a ordem dos ensaios e registre os resultados nas fichas ao final; não há resultados de hardware preenchidos previamente.
\end{abstract}
\begin{IEEEkeywords}
STM32F4-Discovery, WAV PCM, USB Host, FatFs, DMA, áudio embarcado.
\end{IEEEkeywords}}
\maketitle
\IEEEdisplaynontitleabstractindextext
\IEEEpeerreviewmaketitle

\section{Introdução e objetivos}
O firmware de referência demonstra a reprodução de áudio usando o codec da placa STM32F4-Discovery. A variante \texttt{MEDIA\_InFLASH} toca amostras incorporadas no executável; \texttt{MEDIA\_USB\_KEY} atua como USB Host MSC, monta o volume com FatFs e lê um WAV em blocos enquanto o DMA alimenta a interface de áudio. A nota AN3997 descreve o exemplo e sua arquitetura \cite{an3997}.

Ao longo da prática, primeiro confirme o exemplo original sem alterar fontes. Só depois substitua a faixa e experimente a leitura USB. Essa ordem ajuda a separar problemas de configuração do projeto de problemas no formato do áudio. Ao final de cada ensaio, registre o que observou na placa e guarde o log de compilação.

\section{Materiais e ferramentas}
\begin{itemize}
\item Placa STM32F4-Discovery com STM32F407VG, cabo USB de depuração e saída de áudio adequada para conferir o codec;
\item Atollic TrueSTUDIO compatível com o projeto Eclipse fornecido; ST-LINK drivers/utilitários instalados;
\item pacote de firmware \texttt{STM32F4-Discovery\_FW\_V1.1.0} já presente na raiz;
\item repositório do projeto: \url{https://github.com/amarconn/audio-pendrive-stm32};
\item pendrive USB Mass Storage, recomendado FAT32, e computador para prepará-lo;
\item computador com Python 3 e FFmpeg no PATH para converter as faixas; yt-dlp só é necessário se a entrada vier de uma URL.
\end{itemize}

O arquivo \path{audio/smw_course_clear.wav} veio do site \textit{The Mushroom Kingdom}, que reúne efeitos sonoros e trechos musicais de jogos do Mario em formato WAV \cite{tmk}. Ele será usado nos ensaios de conversão.

\section{Pacote V1.1.0 e estrutura do projeto}
O pacote \texttt{STM32F4-Discovery\_FW\_V1.1.0} inclui CMSIS, STM32F4 Standard Peripheral Library, USB Host, FatFs e os utilitários da placa \cite{fw}. Os dois projetos ficam em \path{Project/Audio_playback_and_record/TrueSTUDIO}. Abra primeiro \texttt{MEDIA\_InFLASH} e depois \texttt{MEDIA\_USB\_KEY}. Os arquivos Eclipse usam links relativos para as bibliotecas nas pastas vizinhas; mantenha a árvore do pacote completa. Esta é a versão pedida para a prática e já está disponível junto aos materiais.

\subsection{Importação no Atollic}
Abra o workspace do TrueSTUDIO, escolha \textit{File > Import > General > Existing Projects into Workspace}, selecione como raiz a pasta \texttt{.../TrueSTUDIO} e importe cada projeto. Não marque cópia para workspace: os links Eclipse usam a posição relativa do pacote. Se um projeto aparecer com caminhos quebrados, confirme primeiro que o pacote manteve sua hierarquia e que o workspace não recebeu uma cópia isolada. Selecione o alvo STM32F407VG e a configuração de depuração ST-LINK indicada pelo projeto. O projeto inclui \texttt{.launch}, \texttt{.cproject} e linker script; use-os como configuração inicial e só ajuste o probe/driver do ST-LINK se a instalação local exigir.

Para o teste de referência, faça \textit{Project > Clean}, depois \textit{Build Project}; registre o console e a configuração ativa. Para gravar, conecte a placa pelo conector ST-LINK, selecione \textit{Debug As > Embedded C/C++ Application} (ou a opção equivalente da versão instalada), permita a gravação pelo debugger, e use \textit{Resume}. Os nomes de menus podem variar entre versões do Atollic. Antes de executar, confirme que a configuração é a correta: \texttt{MEDIA\_InFLASH} define a mídia interna; \texttt{MEDIA\_USB\_KEY} define USB Host.

\section{Teste inicial: Flash interna}
Em \texttt{MEDIA\_InFLASH}, a macro da mídia seleciona o caminho Flash em \texttt{main.c}; o programa chama \texttt{WavePlayBack(I2S\_AudioFreq\_48k)}. O áudio de referência está no vetor \texttt{AUDIO\_SAMPLE} de \texttt{src/audio\_sample.c}, em palavras de 16 bits. O código de reprodução obtém o endereço de início e tamanho em \texttt{src/waveplayer.c} (\texttt{AUIDO\_START\_ADDRESS} e \texttt{AUDIO\_FILE\_SZE}) e aciona o codec por \texttt{AudioFlashPlay}. Assim, trocar só o arquivo WAV no disco não muda o conteúdo executado: os bytes precisam ser incorporados ao vetor e os parâmetros coerentes com o formato gerado.

No primeiro ensaio, não altere \emph{nenhum} fonte. Importe, compile, grave e execute a versão original, então confira a saída analógica com fones ou uma caixa apropriada. Anote a configuração de build, mensagens, LEDs e o que ouviu. Se algo falhar, registre a mensagem e o estado dos LEDs antes de modificar qualquer código. Ao final, preencha a linha correspondente na ficha de resultados.

\section{Substituição do áudio na Flash}
O vetor de áudio fica em \texttt{src/audio\_sample.c}; o tamanho e o ponto de início são definidos em \texttt{src/waveplayer.c}. O exemplo inicia o codec a 48 kHz. Para facilitar a troca de faixa, foi preparado um conversor que recebe o WAV, ajusta o áudio para estéreo PCM de 16 bits a 48 kHz e grava o resultado como uma tabela \texttt{uint16\_t}. O cabeçalho de saída imita o formato de 58 bytes usado pela tabela de referência. O script limita a saída a 990.058 bytes para respeitar o tamanho nominal usado pelo exemplo; o espaço restante precisa acomodar o programa.

O arquivo \path{audio/smw_course_clear.wav} foi conferido: é PCM mono, 16 bits, 22.050 Hz, 7,491 s e 330.432 bytes. Ao byte 36 aparece o chunk \texttt{LIST}; por isso, esse arquivo não deve ser copiado diretamente para o pendrive do exemplo USB. Para a Flash, cinco segundos em estéreo 48 kHz ocupam 960.000 bytes de amostras. O vetor de demonstração já gerado está em \path{audio/smw_course_clear_flash.c}. Essa conversão demonstra a preparação dos dados, mas não comprova reprodução na placa.

Para gerar novamente o vetor a partir da raiz do projeto:
\begin{lstlisting}
python scripts/wav_to_flash_array.py audio/smw_course_clear.wav --seconds 5 --output audio/smw_course_clear_flash.c
\end{lstlisting}
O parâmetro \texttt{--seconds 5} seleciona os primeiros cinco segundos, reduzindo a faixa original de 7,491 s para caber no orçamento nominal. O utilitário informa o valor de \texttt{AUDIO\_FILE\_SZE}; para este arquivo, use \texttt{960000u} e mantenha \texttt{AUIDO\_START\_ADDRESS} em 58. A diferença decorre do código ST: o ponteiro avança 58 elementos de 16 bits, enquanto o tamanho enviado ao codec é calculado subtraindo 58 do macro. Substitua \texttt{audio\_sample.c} e ajuste essas constantes apenas em uma cópia de \texttt{MEDIA\_InFLASH}. Em seguida, limpe e recompile o projeto. Se o linker indicar falta de Flash, reduza \texttt{--seconds}; o script limita o áudio, mas o mapa do linker é que confirma o espaço total ocupado.

\section{Conversores de áudio}
Para o pendrive, use \path{scripts/converter_audio_stm32_corrigido.py}. Ele normaliza a faixa com FFmpeg e só grava a saída depois de validar o cabeçalho. Exige Python 3 e FFmpeg no \texttt{PATH}; para origem por URL, também requer yt-dlp. Execute a partir da raiz do projeto:
\begin{lstlisting}
python -m pip install yt-dlp       # somente para URL
python scripts/converter_audio_stm32_corrigido.py audio/smw_course_clear.wav -o audio/audio_usb.wav
python scripts/converter_audio_stm32_corrigido.py "URL" --youtube -o audio/audio_usb.wav
\end{lstlisting}
O conversor USB produz PCM signed little-endian, estéreo, 16 bits e 44.100 Hz; remove metadados e verifica RIFF/WAVE, \texttt{fmt } de 16 bytes e \texttt{data} no offset 36. FFmpeg também é necessário para WAV, pois frequência, canais e chunks variam entre arquivos. Para gerar a tabela da Flash, use \path{scripts/wav_to_flash_array.py}, descrito na seção anterior. O áudio \path{audio/robocop.wav} pode ser usado como entrada do conversor USB.

\section{MEDIA\_USB\_KEY: funcionamento e arquivo esperado}
Em \texttt{main.c}, a variante USB inicializa o botão USER como interrupção e chama \texttt{USBH\_Init} com classe MSC; o laço principal executa \texttt{USBH\_Process}. A sequência do host enumera a memória, monta FatFs e entra na aplicação. \texttt{WAVE\_NAME} em \texttt{inc/main.h} vale \texttt{"0:audio.wav"}; o volume lógico \texttt{0:} é o único drive configurado e o caminho não contém subdiretório, portanto o arquivo deve estar na raiz com esse nome (respeitando a configuração sem suporte a nomes longos).

O \texttt{ffconf.h} incluído habilita escrita (\texttt{\_FS\_READONLY=0}), desabilita LFN (\texttt{\_USE\_LFN=0}), configura um drive (\texttt{\_DRIVES=1}) e define \texttt{\_MAX\_SS=1024} \cite{fatfs}. O driver de disco deve suportar a consulta de tamanho de setor quando o máximo excede 512; esta é uma dependência real da configuração. FAT32 é uma recomendação prática de interoperabilidade para mídia USB; confirme que o pendrive usa uma partição primária compatível. Formate no computador como FAT32 e copie o arquivo final como \texttt{audio.wav} para a raiz, ejetando com segurança.

\subsection{Parser WAV e limitações}
\funcao{WavePlayerStart()} tenta abrir o diretório \texttt{0:/}, depois abre \texttt{WAVE\_NAME}, lê o primeiro bloco e chama a rotina abaixo para validar o cabeçalho:
\begin{lstlisting}
WavePlayer_WaveParsing()
\end{lstlisting}
O parser verifica RIFF no offset 0, WAVE em 8, \texttt{fmt } em 12, formato PCM em 20, canais em 22, taxa em 24 e 16 bits em 34. Ele aceita \texttt{fmt} size diferente de 16 somente em um caso particular de extensão \texttt{fact}; depois procura \texttt{data} no offset derivado. O WAV simplificado de 44 bytes evita essa ramificação e coloca \texttt{data} em 36. Chunks arbitrários como LIST/JUNK não são percorridos genericamente e podem deslocar campos e causar rejeição.

O firmware aceita um ou dois canais, amostra de 16 bits e a taxa que o codec reconhece; o roteiro usa estéreo 44,1 kHz porque é a especificação do conversor. Em caso de dúvida, use essa especificação exata. A implementação mantém \texttt{buffer1} e \texttt{buffer2} de tamanho \texttt{\_MAX\_SS} elementos de 16 bits, lê \texttt{\_MAX\_SS} bytes por chamada FatFs e alterna o buffer consumido pelo DMA com o outro preenchido via \texttt{f\_read}. O tamanho configurado representa 1024 bytes por leitura/DMA; não implica limite de arquivo de 1 MB. \texttt{WaveDataLength} vem do chunk data e é decrementado a cada transferência; arquivos maiores são lidos em várias operações, sujeitos aos limites de FatFs, capacidade do volume e correta leitura final. O código usa \texttt{BytesRead}, mas o fluxo normal não trata cada leitura curta como um caso especial, o que merece atenção para arquivos truncados ou erros de I/O.

Quando uma transferência DMA termina, a rotina abaixo marca o bloco como concluído, atualiza os bytes restantes e permite que o laço carregue o próximo bloco:
\begin{lstlisting}
void EVAL_AUDIO_TransferComplete_CallBack(
    uint32_t pBuffer, uint32_t Size);
\end{lstlisting}
No projeto USB, o comando inicial é \texttt{CMD\_PLAY} (0); depois que o volume é montado, o exemplo tenta abrir e reproduzir o arquivo. A rotina \funcao{EXTI0\_IRQHandler()} alterna \texttt{Command\_index} entre PLAY e \texttt{CMD\_RECORD} (1); ao passar para gravação, para o codec. USER não é um controle play/pause: durante a reprodução, o acionamento muda para gravação e interrompe o áudio; o acionamento seguinte retorna ao comando PLAY. Na variante Flash, o botão atua em pausa e continuação. Se o pendrive for removido durante a reprodução, o firmware chama a rotina de reinicialização do player.

\subsection{Diagnóstico pelos LEDs}
Os significados observáveis no fonte são:
\begin{itemize}
\item LED verde (LED4) aceso indica início/aplicação e aparece no evento de dispositivo anexado;
\item LED vermelho (LED5) aceso/toggle indica ausência de dispositivo/erro em algumas rotas; falha em \texttt{f\_opendir("0:/")} alterna vermelho indefinidamente;
\item falha ao abrir o arquivo WAV acende vermelho e define comando de falha, sem entrar no laço de alternância usado pelo erro de diretório;
\item WAV rejeitado pelo parser alterna vermelho indefinidamente, distinguindo-se da falha de abertura pelo padrão de piscar;
\item LED azul (LED6) alterna durante reprodução e apaga/para durante pausa; LED verde pode indicar término conforme política de repetição.
\end{itemize}
Os comentários e estados estão em \texttt{waveplayer.c}, \texttt{stm32f4xx\_it.c} e \texttt{usbh\_usr.c}; a indicação depende também do timer de LEDs. Confirme no hardware para observar a cadência real.

\section{Procedimento experimental USB}
\begin{enumerate}
\item Importe e compile o projeto original \texttt{MEDIA\_USB\_KEY} sem editar fontes. Grave/execute conforme a seção 3 e registre build e mensagens.
\item Gere um WAV simples com o comando da seção 6. O validador imprime tamanho e duração; corrija a entrada se mostrar erro.
\item Formate o pendrive como FAT32 e copie somente \texttt{audio.wav} para a raiz. No projeto confirme \texttt{WAVE\_NAME "0:audio.wav"}.
\item Conecte o pendrive à porta USB Host prevista pelo hardware, ligue a placa e aguarde o reconhecimento e a montagem do volume. Com o arquivo correto, o firmware inicia a reprodução pelo comando PLAY.
\item Durante a reprodução, pressione USER uma vez: o código muda para RECORD, interrompe o áudio e inicia o fluxo de gravação. Pressione novamente para retornar a PLAY. Registre o comportamento; não deixe gravações seguirem por tempo desnecessário.
\item Repita com WAV deliberadamente incompatível apenas se for seguro e anote a diferença entre arquivo ausente, diretório não montado e parser rejeitando o cabeçalho.
\end{enumerate}
Execute estes passos com a placa e registre as observações no caderno de laboratório. A figura a seguir resume o caminho dos dados; ela é um esquema explicativo, não uma captura da bancada.

\begin{figure}[!ht]
\centering
\begin{tikzpicture}[node distance=0.25cm, every node/.style={draw,rounded corners,align=center,font=\footnotesize,minimum height=0.55cm,text width=3.6cm},arr/.style={->,thick}]
\node (usb) {Pendrive: WAV PCM};
\node[below=of usb] (host) {USB Host: MSC};
\node[below=of host] (fs) {FatFs: leitura de 1024 bytes};
\node[below=of fs] (dma) {Buffers alternados e DMA};
\node[below=of dma] (codec) {Codec e saída de áudio};
\draw[arr] (usb)--(host); \draw[arr] (host)--(fs); \draw[arr] (fs)--(dma); \draw[arr] (dma)--(codec);
\end{tikzpicture}
\caption{Fluxo conceitual de reprodução USB conforme o código.}
\label{fig:fluxo}
\end{figure}

\section{Conclusão}
Ao terminar, compare os três ensaios: Flash original, Flash com nova faixa e leitura pelo pendrive. Use os registros para localizar em que etapa surgiu cada diferença. Guarde o log do build e os valores de conversão junto às anotações para que outro aluno consiga repetir o procedimento.

\begin{thebibliography}{3}
\bibitem{an3997} STMicroelectronics, \textit{AN3997: Audio playback and recording using the STM32F4DISCOVERY}, nota de aplicação fornecida no repositório.
\bibitem{fw} STMicroelectronics, \textit{STM32F4-Discovery Firmware Package V1.1.0}, pacote de firmware fornecido no repositório; exemplo Audio playback and recording.
\bibitem{fatfs} Elm-Chan, \textit{FatFs -- Generic FAT Filesystem Module}, versão/configuração incluída no pacote de firmware analisado.
\bibitem{tmk} The Mushroom Kingdom, \textit{Sound Clips: voice clips, sound effects, and short musical themes in WAV format}, \url{https://themushroomkingdom.net/wav.shtml}.
\end{thebibliography}
\end{document}
